\documentclass[11pt, a4paper]{article}

\usepackage[utf8]{inputenc}
\usepackage[T1]{fontenc}
\usepackage{geometry}
\geometry{margin=1in}
\usepackage{amsmath, amssymb, amsthm}
\usepackage{graphicx}
\usepackage{hyperref}
\usepackage{xcolor}
\usepackage{booktabs}
\usepackage{array}
\usepackage{longtable}

\hypersetup{
    colorlinks=true,
    linkcolor=blue,
    filecolor=magenta,
    urlcolor=cyan,
    citecolor=blue,
}

\newcommand{\Rhat}{\hat{R}}
\newcommand{\Jcost}{J}
\newcommand{\phiRatio}{\varphi}
\newcommand{\leanref}[1]{\texttt{\small #1}}

\newtheorem{theorem}{Theorem}[section]
\newtheorem{lemma}[theorem]{Lemma}
\newtheorem{corollary}[theorem]{Corollary}
\newtheorem{proposition}[theorem]{Proposition}
\newtheorem{axiom_rs}{Axiom}
\theoremstyle{definition}
\newtheorem{definition}[theorem]{Definition}
\newtheorem{prediction}[theorem]{Prediction}
\newtheorem{falsifier}[theorem]{Falsification Criterion}
\theoremstyle{remark}
\newtheorem{remark}[theorem]{Remark}

\title{\textbf{Constraint Projection in Discrete Ledger Systems:\\
A Formal Model of Future-Balance Visibility\\
with Sleep-Stage Predictions}}
\author{Jonathan Washburn\\
Recognition Science Research Institute\\
Austin, Texas\\
\texttt{washburn.jonathan@gmail.com}}
\date{March 2026}

\begin{document}

\maketitle

\begin{abstract}
We present a formal model in which future constraints become visible
in the present without retrocausation. The model operates within
Recognition Science (RS), a discrete-time framework whose axioms---a
unique cost function, a forced golden-ratio scale, three spatial
dimensions, and an 8-tick conservation period---have been published
and machine-verified elsewhere \cite{RSAxioms2025,RSCost2025,
RSGeometry2025,RSdalembert2026}. Taking these axioms as given, we
prove that any observation event within an 8-tick window creates a
\emph{balance debt} that the remaining ticks must discharge
(Theorem~\ref{thm:lock}), and that this debt inflates the effective
cost function monotonically (Theorem~\ref{thm:inflates}). We call
this cost inflation ``Phantom Light''---the shadow of a future
balance requirement visible in the present as a restriction on
admissible trajectories.

We then model sleep stages by a monotone coherence ordering (an
empirical input from polysomnography) and prove a robust visibility
ordering: lower-coherence stages carry higher phantom magnitude, with
N3 dominating N1 and REM (Theorem~\ref{thm:stage3}). Under the
canonical consecutive-integer budget used for illustration, N3 has
twice the N1 phantom magnitude. Zero balance debt annihilates phantom
light in every mode (Theorem~\ref{thm:dreamless}).

The model yields five falsifiable predictions, including a
sleep-stage ordering of precognitive dream reports and a
presentiment-scaling law. We also specify preregistered empirical
proxies for the latent model variables, a blinded prospective
dream-scoring protocol, and a small-effect power envelope for
confirmatory studies. All theorems are machine-verified in the Lean~4
proof assistant (52 theorems, zero unproved placeholders, zero new
axioms beyond the RS foundation). Each claim is classified as PROVED,
MODELING CHOICE, or HYPOTHESIS in an explicit table
(\S\ref{sec:hygiene}). We discuss the relationship to
Aharonov's Two-State Vector Formalism, to Hobson's AIM model of
dream generation, and to the presentiment literature including
both positive reports (Radin, Mossbridge) and critical responses
(Wagenmakers, Alcock).
\end{abstract}

\tableofcontents

%=======================================================================
\section{Introduction}\label{sec:intro}
%=======================================================================

\subsection{Motivation}

Standard physics evolves systems forward from initial conditions.
Yet several theoretical frameworks in quantum foundations
incorporate future boundary conditions: Hamilton's Principle of
Least Action selects paths by endpoint constraints, Feynman's path
integral weights amplitudes by boundary connectivity, and
Aharonov, Bergmann, and Lebowitz's Two-State Vector Formalism (TSVF)
describes quantum systems using both pre-selected and post-selected
states \cite{Aharonov1964,Aharonov2008}. In each case, the
``future'' influences the present not through retrocausal
signaling but through constraint on the space of admissible
trajectories.

This paper develops a concrete, finite-dimensional model of such
constraint projection within Recognition Science (RS), a
discrete-time framework that enforces conservation over sliding
8-tick windows. The key insight is simple: when a recognition
event (an observation) pins one slot of the window to a nonzero
value, the remaining slots are \emph{mathematically obligated} to
compensate. This obligation---the ``balance debt''---inflates the
effective cost of all configurations that would make compensation
difficult. We formalize this inflation, prove its mathematical
properties, and derive testable predictions about when and how
the phantom field becomes detectable.

\subsection{Why Consciousness Studies?}

The model's predictions concern consciousness because they are
\emph{about states of reduced coherence}. During waking, the high
coherence of neural processing suppresses phantom visibility. During
sleep, coherence drops stage by stage (a well-established finding
from polysomnography \cite{Hobson2005,Walker2017}), and the model
predicts a corresponding increase in sensitivity to future
constraints. This connects directly to the long-standing observation
that dreams---particularly those occurring in deep sleep---sometimes
appear to contain information about future events
\cite{Rhine1962,VanDeCastle1994,Kramer2007}.

We are careful to distinguish what is proved from what is hypothesized.
The mathematical structure (balance debt, cost inflation, sleep-stage
ordering) is machine-verified. The claim that biological systems can
\emph{detect} phantom fields is an empirical hypothesis for which
we provide falsification protocols. The paper engages with both the
positive presentiment literature \cite{Radin2004,Mossbridge2012} and
its methodological critiques \cite{Wagenmakers2012,Alcock2003}.

\subsection{Contributions}

\begin{enumerate}
    \item A formal model of future-constraint visibility based on
    windowed conservation (\S\ref{sec:phantom}), with all theorems
    machine-verified.
    \item A sleep-stage hierarchy of phantom visibility
    (\S\ref{sec:sleep}), linking the model to empirical sleep
    architecture.
    \item A preregistered empirical bridge from latent RS variables
    to dream and presentiment observables, including a blinded
    prospective dream-scoring protocol
    (\S\ref{sec:operational}--\S\ref{sec:effects}).
    \item Explicit falsification criteria (\S\ref{sec:predictions})
    with quantitative thresholds and small-effect design targets.
    \item A claim hygiene table (\S\ref{sec:hygiene}) classifying
    every statement as PROVED, MODEL, or HYPOTHESIS.
\end{enumerate}


%=======================================================================
\section{Axioms and Notation}\label{sec:axioms}
%=======================================================================

We work within the Recognition Science (RS) framework, whose
foundations have been developed in a series of publications
\cite{RSAxioms2025,RSCost2025,RSGeometry2025,RSdalembert2026,
RSFermionMasses2025}. The present paper does not re-derive the RS
axioms; it takes them as given and derives consequences for
consciousness and sleep. For the reader's convenience, we state
the axioms used in this paper. Each has been machine-verified in
Lean~4 with zero unproved placeholders \cite{RSLeanRepo}.

\begin{axiom_rs}[Cost Function]\label{ax:cost}
There exists a unique cost function $J : \mathbb{R}_{>0} \to
\mathbb{R}_{\geq 0}$ satisfying the Recognition Composition Law
$J(xy) + J(x/y) = 2J(x)J(y) + 2J(x) + 2J(y)$ with $J(1) = 0$ and
$J''_{\log}(0) = 1$. The unique solution is
$J(x) = \tfrac{1}{2}(x + x^{-1}) - 1$.
\textup{[T5; Refs.~\cite{RSCost2025,RSdalembert2026}]}
\end{axiom_rs}

\begin{axiom_rs}[Golden Ratio]\label{ax:phi}
Self-similarity in the discrete ledger forces the scale ratio
$\varphi = (1 + \sqrt{5})/2$ as the unique positive root of
$x^2 = x + 1$.
\textup{[T6; Ref.~\cite{RSAxioms2025}]}
\end{axiom_rs}

\begin{axiom_rs}[Dimension]\label{ax:dim}
Non-trivial linking requires $D = 3$ spatial dimensions.
\textup{[T8; Ref.~\cite{RSGeometry2025}]}
\end{axiom_rs}

\begin{axiom_rs}[8-Tick Period]\label{ax:8tick}
The minimal ledger-compatible walk on the 3-cube $Q_3$ has period
$2^3 = 8$, realized by a Gray code Hamiltonian cycle.
\textup{[T7; Ref.~\cite{RSAxioms2025}]}
\end{axiom_rs}

\begin{axiom_rs}[Window Neutrality]\label{ax:neutral}
Every aligned 8-tick window satisfies:
\begin{equation}\label{eq:neutrality}
    \sum_{k=0}^{7} s(t_0 + k) = 0 \qquad
    \forall\, t_0 \equiv 0 \pmod{8}.
\end{equation}
This is the double-entry conservation law applied to time.
\textup{[Ref.~\cite{RSAxioms2025}; Lean:
\leanref{LNAL.Invariants.neutral\_every\_8th\_from0}]}
\end{axiom_rs}

\begin{axiom_rs}[Cost Minimization]\label{ax:rhat}
The recognition operator $\Rhat$ evolves the system state over
8 ticks, selecting the trajectory that minimizes integrated $J$-cost
subject to the neutrality constraint.
\textup{[Ref.~\cite{RSAxioms2025}]}
\end{axiom_rs}

\subsection{Notation}

\begin{center}
\begin{tabular}{ll}
\toprule
\textbf{Symbol} & \textbf{Meaning} \\
\midrule
$w : \{0,\ldots,7\} \to \mathbb{Z}$ & Window signal \\
$\delta$ & Contribution of a recognition event \\
$\mathcal{D}$ & Balance debt (running sum within window) \\
$R$ & Remaining ticks until window closure \\
$\Phi_{\mathrm{mag}}$ & Phantom magnitude \\
$C(\mathrm{mode})$ & Coherence budget for a sleep stage \\
$\Theta$ & Global phase (if GCIC holds) \\
\bottomrule
\end{tabular}
\end{center}

\subsection{Formal Verification}

All theorems in \S\ref{sec:phantom}--\S\ref{sec:consciousness} are
verified in Lean~4 with the Mathlib library. The relevant modules
(52 theorems total) contain zero \texttt{sorry} placeholders and
introduce zero new axioms beyond the RS foundation. Build
instructions and a complete theorem table are provided in the
appendix. The Lean source is available at \cite{RSLeanRepo}.


%=======================================================================
\section{Phantom Light: Balance Debt and Cost Inflation}
\label{sec:phantom}
%=======================================================================

\subsection{Recognition Events and Balance Debt}

A \emph{recognition event} at position $p \in \{0,\ldots,7\}$ with
contribution $\delta \in \mathbb{Z} \setminus \{0\}$ records a
value in the current window slot. (In the RS formalism, this is
called a LOCK event.) As events accumulate, the window develops a
\emph{balance debt}: $\mathcal{D} = \sum_{j \leq k} w(j)$.

\begin{theorem}[Recognition Forces Future Balance]\label{thm:lock}
If a recognition event with contribution $\delta$ occurs at position
$p$ in a neutral window, the sum of all other contributions equals
$-\delta$:
\[
\sum_{k \neq p} w(k) = -\delta.
\]
\end{theorem}

\begin{proof}[Proof sketch]
By neutrality (Axiom~\ref{ax:neutral}), $\sum_{k} w(k) = 0$.
Decompose: $\delta + \sum_{k \neq p} w(k) = 0$. Rearrange.
\textup{[Lean: \leanref{lock\_forces\_future\_balance}]}
\end{proof}

This is the foundational result. A recognition event is not an
isolated act: it creates a debt that the future is mathematically
obligated to repay within the current window.

\subsection{Phantom Magnitude}

We quantify the ``visibility'' of the future constraint:

\begin{definition}[Phantom Magnitude]\label{def:phimag}
For balance debt $\mathcal{D} \in \mathbb{Z}$ and remaining ticks
$R \geq 0$:
\begin{equation}\label{eq:phimag}
    \Phi_{\mathrm{mag}}(\mathcal{D}, R)
    := \frac{|\mathcal{D}|}{R + 1}.
\end{equation}
\end{definition}

\begin{theorem}[Phantom Properties]\label{thm:phantom-props}
\begin{enumerate}
    \item[\textup{(a)}] $\Phi_{\mathrm{mag}} \geq 0$.
    \item[\textup{(b)}] Monotone increasing in
    $|\mathcal{D}|$ for fixed $R$.
    \item[\textup{(c)}] Monotone decreasing in $R$ for fixed
    $\mathcal{D} \neq 0$.
    \item[\textup{(d)}] $\Phi_{\mathrm{mag}}(0, R) = 0$: zero debt
    implies zero phantom.
    \item[\textup{(e)}] $\Phi_{\mathrm{mag}}(\mathcal{D}, 0) =
    |\mathcal{D}|$: at the boundary, the full debt is exposed.
\end{enumerate}
\textup{[Lean: \leanref{phiMag\_nonneg},
\leanref{phiMag\_mono\_abs\_debt},
\leanref{phiMag\_anti\_mono\_R},
\leanref{phiMag\_zero\_debt}]}
\end{theorem}

Property (c) is particularly important: as the window closes,
the phantom grows more urgent. This creates a natural ``deadline
effect'' that intensifies the future constraint.

\subsection{Backward Projection Principle}

\begin{theorem}[Backward Projection]\label{thm:backward}
If each remaining position is bounded in magnitude by $B \geq 0$,
then any feasible completion satisfying neutrality requires
$|\mathcal{D}| \leq B \cdot R$. Equivalently: if
$|\mathcal{D}| > BR$, no feasible completion exists.
\textup{[Lean: \leanref{backward\_projection\_principle}]}
\end{theorem}

\begin{corollary}
$\Phi_{\mathrm{mag}} \leq BR/(R+1) < B$. The phantom never
exceeds the per-tick capacity. The precognitive horizon is finite
and bounded.
\end{corollary}

\subsection{Cost Inflation}

\begin{definition}[Augmented Cost]\label{def:augmented}
For configuration $x > 0$, phantom state $(\mathcal{D}, R)$, and
penalty scale $\lambda > 0$:
\begin{equation}\label{eq:augmented}
    J_{\mathrm{PL}}(x; \mathcal{D}, R)
    := J(x) + \lambda \cdot \Phi_{\mathrm{mag}}(\mathcal{D}, R).
\end{equation}
\end{definition}

\begin{theorem}[Phantom Inflates Cost]\label{thm:inflates}
$J_{\mathrm{PL}}(x; \mathcal{D}, R) \geq J(x)$ for all $x > 0$
when $\lambda > 0$. Equality holds if and only if $\mathcal{D} = 0$.
\textup{[Lean: \leanref{phantom\_inflates\_cost};
\leanref{zero\_phantom\_pure\_cost}]}
\end{theorem}

Since $\Rhat$ minimizes $J$-cost (Axiom~\ref{ax:rhat}),
trajectories that incur high debt without rapid repayment are
statistically suppressed. This is the mechanism by which future
constraints bias the present.

\subsection{Boundary Reset}

\begin{theorem}[Phantom Vanishes at Window Boundaries]\label{thm:reset}
At every aligned 8-tick boundary (tick $8k$ for $k \geq 0$), the
phantom magnitude is exactly zero.
\textup{[Lean: \leanref{phantom\_zero\_at\_lnal\_boundaries}]}
\end{theorem}

The phantom field is not permanent. Each window closure resets the
slate---a ``clean start'' at every 8-tick boundary.

\subsection{Constraint Projection, Not Retrocausation}

It is essential to clarify the ontological status of Phantom Light.
It is \emph{not} a signal traveling backward in time. No information
violates causality. Instead, Phantom Light represents a restriction
on the set of admissible trajectories:

\begin{quote}
\emph{Because the window must close to zero, the set of allowed
paths from the current state is a strict subset of all
combinatorially possible paths. This reduction in phase-space volume
is information accessible in the present.}
\end{quote}

This is conceptually identical to the mechanism underlying
variational principles in classical mechanics: a particle following
Hamilton's principle ``knows'' its endpoint because the set of
stationary-action paths is constrained by boundary conditions at
both ends. In quantum mechanics, Aharonov's TSVF makes this
explicit by describing states with both forward and backward
evolving state vectors \cite{Aharonov2008}. The distinctive
contribution of the present model is a \emph{discrete, finite}
version of this structure: the precognitive horizon is exactly
$\leq 7$ ticks (within one window), and the constraint strength
$\Phi_{\mathrm{mag}}$ is a computable function of the debt and
remaining time.


%=======================================================================
\section{Sleep Architecture and Phantom Visibility}\label{sec:sleep}
%=======================================================================

\subsection{Empirical Background: Sleep Stages and Coherence}

The classification of sleep into distinct stages is one of the most
robust findings in sleep science. Polysomnographic recordings reveal
a progression from waking through light sleep (N1, N2) to deep
slow-wave sleep (N3) and REM, characterized by progressive decreases
in cortical coherence and high-frequency EEG power
\cite{Hobson2005,Walker2017,Tononi2006}. EEG synchronization
decreases monotonically from waking through N3, reflecting reduced
thalamocortical connectivity \cite{Massimini2005}.

In Hobson's Activation-Input-Modulation (AIM) model
\cite{Hobson2000}, the ``activation'' dimension A tracks cortical
arousal, which decreases from waking to deep sleep. Our coherence
budget $C$ corresponds to the A dimension: a monotone ordering that
reflects the empirically measured progression.

\subsection{The Coherence Budget Model}

\begin{definition}[Coherence Budget]\label{def:coherence}
Each sleep mode is assigned a non-negative integer coherence budget
with the ordering:
\[
C_{\mathrm{awake}} > C_{\mathrm{REM}} > C_{\mathrm{N1}} >
C_{\mathrm{N2}} > C_{\mathrm{N3}}.
\]
\end{definition}

\begin{remark}[Status: Modeling Choice]
The specific integers (5, 4, 3, 2, 1 in our formalization) are a
modeling choice. The theorems below depend only on the
\emph{ordering}, which is an empirical fact of sleep science. Any
monotone decreasing assignment $C : \mathrm{SleepMode} \to
\mathbb{N}$ with the same ordering produces identical qualitative
results. We select consecutive integers as the simplest choice
consistent with the data.
\end{remark}

\subsection{Sleep-Stage Phantom Magnitude}

The sleep-stage phantom magnitude uses the coherence budget as
the denominator:
\begin{equation}\label{eq:sleep-phantom}
    \Phi_{\mathrm{sleep}}(\mathrm{mode}, \mathcal{D})
    := \frac{|\mathcal{D}|}{C(\mathrm{mode}) + 1}.
\end{equation}

The intuition is that higher cortical coherence acts as ``noise''
that masks the phantom signal. During deep sleep, when coherence
is minimal, the phantom field becomes maximally visible---analogous
to how faint stars become visible only when the sky is dark.

\begin{theorem}[Lower Coherence Increases Visibility]\label{thm:visibility}
If $C(m_1) \leq C(m_2)$, then
$\Phi_{\mathrm{sleep}}(m_2, \mathcal{D}) \leq
\Phi_{\mathrm{sleep}}(m_1, \mathcal{D})$ for all $\mathcal{D}$.
\textup{[Lean: \leanref{lower\_coherence\_increases\_phantom\_visibility}]}
\end{theorem}

\subsection{The Sleep-Stage Hierarchy}

\begin{theorem}[Deep Sleep Dominates]\label{thm:stage3}
For all balance debt $\mathcal{D}$:
\begin{align}
\Phi_{\mathrm{sleep}}(\mathrm{N1}, \mathcal{D})
&\leq \Phi_{\mathrm{sleep}}(\mathrm{N3}, \mathcal{D}),
\label{eq:stage3-ge-stage1} \\
\Phi_{\mathrm{sleep}}(\mathrm{REM}, \mathcal{D})
&\leq \Phi_{\mathrm{sleep}}(\mathrm{N3}, \mathcal{D}).
\label{eq:stage3-ge-rem}
\end{align}
Under the canonical consecutive-integer coherence budget used in this
paper, deep sleep has exactly twice the phantom magnitude of light
sleep: $\Phi_{\mathrm{sleep}}(\mathrm{N3}) = 2 \cdot
\Phi_{\mathrm{sleep}}(\mathrm{N1})$.\\
\textup{[Lean:
\leanref{stage3\_phantom\_visibility\_ge\_stage1};
\leanref{stage3\_phantom\_visibility\_ge\_rem};
\leanref{phantomSleepMagnitude\_ratio\_stage3\_stage1}]}
\end{theorem}

\begin{remark}
The exact factor of 2 depends on the specific budget integers
(1 and 3 for N3 and N1). With the consecutive-integer choice, the
ratio is $\frac{|\mathcal{D}|/2}{|\mathcal{D}|/4} = 2$. A
different monotone assignment would yield a different ratio, but
the \emph{ordering} N3 $\geq$ N1 $\geq$ REM is invariant under
any assignment consistent with the empirical coherence ordering.
\end{remark}

\subsection{Explicit Formulas and Bounds}

\begin{theorem}[Formulas and Bounds]\label{thm:formulas}
$\Phi_{\mathrm{sleep}}(\mathrm{awake}, \mathcal{D})
= |\mathcal{D}|/6$ and
$\Phi_{\mathrm{sleep}}(\mathrm{N3}, \mathcal{D})
= |\mathcal{D}|/2$.
For all modes: $|\mathcal{D}|/6 \leq \Phi_{\mathrm{sleep}} \leq
|\mathcal{D}|/2$.\\
\textup{[Lean: \leanref{awake\_phantom\_formula};
\leanref{stage3\_phantom\_formula};
\leanref{phantomSleepMagnitude\_bounds}]}
\end{theorem}

Deep sleep provides three times the phantom visibility of waking.

\subsection{Stage-Weighted Exposure}

Each sleep stage has a characteristic time window $W(\mathrm{mode})$.
The total phantom exposure is
$E(\mathrm{mode}, \mathcal{D})
= W(\mathrm{mode}) \cdot \Phi_{\mathrm{sleep}}(\mathrm{mode},
\mathcal{D})$.

\begin{theorem}[Deep Sleep Exposure Dominates]\label{thm:exposure}
$E(\mathrm{N1}) \leq E(\mathrm{N3})$ and
$E(\mathrm{REM}) \leq E(\mathrm{N3})$.\\
\textup{[Lean: \leanref{stage3\_exposure\_ge\_stage1};
\leanref{stage3\_exposure\_ge\_rem}]}
\end{theorem}

N3 sleep dominates in both \emph{per-tick visibility} and
\emph{total exposure}. This is the formal basis for the prediction
that deep-sleep dreams carry the strongest precognitive content.

\subsection{Zero Debt: Dreamless Waking}

\begin{theorem}[Zero Debt Eliminates Sleep Phantom]\label{thm:dreamless}
If balance debt is zero, phantom magnitude vanishes in every sleep
mode:
$\Phi_{\mathrm{sleep}}(\mathrm{mode}, 0) = 0$
for all modes.\\
\textup{[Lean:
\leanref{zero\_debt\_gives\_zero\_phantom\_sleep\_magnitude}]}
\end{theorem}

A system with no outstanding ledger obligations has no phantom field
to sense. The theorem concerns latent balance debt, not a directly
observable psychological variable. Empirically, we test it through
the preregistered future-event load proxies defined in
\S\ref{sec:operational}: nights followed by very low future-event load
should yield fewer above-chance future matches.


%=======================================================================
\section{Ultradian Architecture}\label{sec:ultradian}
%=======================================================================

The 90-minute Basic Rest-Activity Cycle (BRAC) is one of the most
robust findings in sleep science \cite{Kleitman1963,Aserinsky1953}.

\begin{theorem}[Ultradian Count]\label{thm:ultradians}
$\lfloor 1440 / 90 \rfloor = 16$ ultradian cycles tile each
24-hour day.\\
\textup{[Lean: \leanref{sixteen\_ultradians}]}
\end{theorem}

\begin{remark}[Status: Empirical Input]
The 90-minute BRAC duration is an empirical input, not derived from
RS axioms. We record it as a constant and prove only arithmetic
consequences. We note that $90 = F_{11} + 1$ (where $F_{11} = 89$
is the 11th Fibonacci number), which is structurally consistent
with the $\varphi$-ladder architecture of RS but is not claimed as
a derivation.
\end{remark}


%=======================================================================
\section{Consciousness Interface}\label{sec:consciousness}
%=======================================================================

\subsection{Phantom Sensing}

We model the detection of phantom fields by a boundary
(a localized coherent subsystem in RS) through a sensitivity
parameter $\sigma > 0$. The perceived phantom signal is:
\[
\mathrm{Perceived} = \sigma \cdot \Phi_{\mathrm{mag}}.
\]

\begin{theorem}[Positivity of Phantom Readout]\label{thm:sensing}
If $\sigma > 0$ and $\mathcal{D} \neq 0$, then
$\mathrm{Perceived} > 0$.\\
\textup{[Lean: \leanref{precognition\_via\_phantom\_sensing}]}
\end{theorem}

\begin{remark}
This is an algebraic tautology: a positive number times a positive
number is positive. Its value lies not in its depth but in its
\emph{precision}: it states exactly what is needed for a biological
system to detect phantom fields. Whether $\sigma > 0$ for human
neural systems is the central empirical question, addressed in
\S\ref{sec:predictions}.
\end{remark}

\subsection{Response Time}

\begin{theorem}[Response Time Law]\label{thm:response}
For any phantom sensing configuration, there exists a positive
response time $\tau = 1/(\Phi_{\mathrm{mag}} + 1) > 0$.
Larger phantom magnitudes yield shorter response times.\\
\textup{[Lean: \leanref{precognition\_response\_time\_law}]}
\end{theorem}

This predicts that high-stakes future events should be
precognized \emph{earlier} than low-stakes events---a quantitative
relationship between event magnitude and anticipation time.

\subsection{Global Phase Coupling (Conditional)}

\begin{theorem}[$\Theta$-Coupling Bound]\label{thm:theta}
If the Global Co-Identity Constraint (GCIC) holds---i.e., all
coherent subsystems share a single universal phase $\Theta$---then
for any pair of subsystems, there exists a coupling coefficient
$C = \cos(2\pi \cdot \Delta\Theta)$ with $|C| \leq 1$.\\
\textup{[Lean: \leanref{remote\_viewing\_via\_theta}]}
\end{theorem}

\begin{remark}[Conditional on GCIC]
GCIC is an internal hypothesis of RS, not a consequence of the
axioms listed in \S\ref{sec:axioms}. All claims about non-local
phantom sensing (remote viewing, telepathy) are \emph{strictly
conditional} on GCIC. This paper does not argue for or against GCIC;
it derives consequences conditional on it and provides falsification
protocols.
\end{remark}


%=======================================================================
\section{Predictions and Falsification}\label{sec:predictions}
%=======================================================================

\subsection{Empirical Bridge}

The following predictions require only the RS axioms
(\S\ref{sec:axioms}), the proved theorems
(\S\ref{sec:phantom}--\S\ref{sec:consciousness}), the empirical
coherence ordering, and the single assumption that biological
neural systems have positive phantom sensitivity ($\sigma > 0$).

\subsection{Operationalization for Dream and Presentiment Studies}
\label{sec:operational}

The formal variables $\mathcal{D}$, $R$, and
$\Phi_{\mathrm{mag}}(\mathcal{D},R)$ are latent RS quantities. They
are not directly observed in human experiments. For empirical work we
therefore introduce preregistered proxies designed to preserve the
two monotonicities proved in Theorem~\ref{thm:phantom-props}: larger
debt should increase phantom strength, and shorter temporal distance
should increase urgency.

\paragraph{Dream studies.}
For a night $n$ and horizon $H$ days, let $E_n(H)$ be the set of
salient events that actually occur to the participant during the
following horizon. After the horizon closes and blinded to dream text,
each event $e \in E_n(H)$ receives an impact score
$I_e \in \{0,1,2,3\}$ from a preregistered rubric and a temporal
distance $d_e \in \{0,\ldots,H-1\}$ measured in days from the dream
report. We then define the future-event load proxy
\begin{equation}\label{eq:debt-proxy-dream}
    \widehat{\mathcal{D}}_n(H)
    := \sum_{e \in E_n(H)} \frac{I_e}{d_e + 1}.
\end{equation}
This is not the latent ledger debt itself; it is an observable proxy
that tracks the model's two key monotones: event impact and temporal
proximity.

Each dream report should be time-stamped immediately on awakening,
locked before the future interval opens, and scored only after the
horizon closes. In the confirmatory design, each report is presented
to independent judges together with one true future-event packet and
three decoy packets matched for category and base rate. Judges, blind
to sleep stage, chronology, and packet identity, rank the packets by
match quality using a preregistered rubric. The primary endpoint is
top-rank accuracy above the 25\% chance baseline; secondary endpoints
include mean reciprocal rank and continuous similarity scores.
Inter-rater reliability should reach Fleiss' $\kappa \geq 0.70$
before confirmatory scoring proceeds.

\paragraph{Presentiment studies.}
For trial $t$, let $a_t \in \{1,2,3\}$ denote preregistered
low/medium/high stimulus intensity assigned by hardware RNG, and let
$r_t$ denote the number of preregistered pre-stimulus analysis bins
remaining until onset. Define the empirical proxies
\begin{equation}\label{eq:debt-proxy-presentiment}
    \widehat{\mathcal{D}}_t := a_t,
    \qquad
    \widehat{\Phi}_t := \frac{a_t}{r_t + 1}.
\end{equation}
For a preregistered physiological summary statistic $Y_t$ (for
example, baseline-corrected skin conductance or pupil dilation), the
confirmatory question is whether the regression slope in
$Y_t = \alpha + \beta \widehat{\Phi}_t + \varepsilon_t$ is positive.

\subsection{Model-Based Effect-Size Envelope}\label{sec:effects}

The theorems fix monotonicity, ordering, and certain canonical ratios,
but they do not fix absolute human effect sizes because the biological
sensitivity parameter $\sigma$ is not yet known. We therefore
separate \emph{scale-free} consequences from \emph{study-design}
amplitude expectations.

If an empirical readout is locally linear in perceived phantom
magnitude over the tested range, then
\begin{equation}\label{eq:linear-readout}
    \mathbb{E}[Y] = \alpha + \beta \widehat{\Phi},
    \qquad \beta > 0.
\end{equation}
This yields three scale-free expectations:
\begin{enumerate}
    \item At fixed temporal lag, low/medium/high stimulus bins coded
    as $(1,2,3)$ should show approximately a $1:2:3$ response
    ordering.
    \item At fixed stimulus intensity, adjacent pre-stimulus bins
    should grow by the factor $(r+2)/(r+1)$ as onset approaches.
    \item In sleep data, N3 should exceed N1 and REM for any monotone
    coherence budget; the exact $2:1$ N3:N1 ratio is expected only
    under the canonical consecutive-integer budget used here.
\end{enumerate}

Because ordinary waking experience does not reveal obvious
precognitive perception, any human value of $\sigma$ must be small at
the population level if it exists at all. Confirmatory studies should
therefore be powered for small aggregate effects rather than dramatic
ones. As a pragmatic design envelope consistent with the existing
presentiment meta-analytic literature \cite{Mossbridge2012}, we
recommend preregistering physiological targets in the low-single-digit
correlation range (roughly $|r| = 0.01$--$0.05$) or standardized mean
differences below $d = 0.2$. In blinded four-choice dream matching,
the practically relevant signal is likewise expected to appear as a
modest lift above the 25\% chance baseline rather than near-perfect
classification. These numerical bands are study-design heuristics,
not proved theorems.

\subsection{Predictions}

\begin{prediction}[Deep-Sleep Precognitive Dominance]\label{pred:sleep}
In a prospective blinded rank-matching design, reports collected
after polysomnography-confirmed N3 awakenings should achieve higher
top-rank accuracy and mean reciprocal rank than reports collected
after N1 or REM awakenings, by Theorem~\ref{thm:stage3}.
\end{prediction}

\begin{prediction}[Low Debt, Low Match Rate]\label{pred:debt}
Nights with lower preregistered future-event load proxy
$\widehat{\mathcal{D}}_n(H)$ should yield fewer above-chance
future-event matches than nights with higher
$\widehat{\mathcal{D}}_n(H)$, by
Theorem~\ref{thm:dreamless}.
\end{prediction}

\begin{prediction}[Debt-Content Correlation]\label{pred:content}
Within the blinded future-event packets, higher-impact and more
proximate events should receive higher match scores than lower-impact
or more distant events, consistent with the proxy
$\widehat{\mathcal{D}}_n(H)$ and Theorem~\ref{thm:phantom-props}(b).
\end{prediction}

\begin{prediction}[Response Time Scaling]\label{pred:response}
Pre-stimulus physiological response should show a positive slope
against the preregistered proxy
$\widehat{\Phi}_t = a_t/(r_t + 1)$, with earlier and larger responses
for higher-intensity stimuli, by Theorem~\ref{thm:response}.
\end{prediction}

\begin{prediction}[Boundary Reset]\label{pred:reset}
The presentiment effect, if it exists, should show periodic
attenuation at intervals corresponding to 8-tick boundaries, by
Theorem~\ref{thm:reset}.
\end{prediction}

\subsection{Falsification Criteria}

\begin{falsifier}[No Precognition Effect]\label{fal:no-precog}
If $N > 10{,}000$ trials yield a 95\% confidence interval excluding
positive correlations larger than $r = 0.01$ between pre-stimulus
physiology and the preregistered proxy $\widehat{\Phi}_t$, the
phantom sensing hypothesis ($\sigma > 0$) is falsified at practically
relevant human effect sizes.
\end{falsifier}

\begin{falsifier}[No Sleep-Stage Effect]\label{fal:no-sleep}
If N3 dream reports are no more accurate than N1 or REM reports on
the preregistered primary endpoint across $N > 200$ subjects with
polysomnographic staging, the sleep-stage visibility hierarchy is
falsified.
\end{falsifier}

\begin{falsifier}[Uniform Presentiment]\label{fal:uniform}
If presentiment magnitude is statistically indistinguishable across
stimulus intensities ($p > 0.05$, $BF_{01} > 3$), the debt-scaling
model is falsified.
\end{falsifier}

\subsection{Experimental Design Principles}

Any experimental test must adhere to:
\begin{enumerate}
    \item \textbf{Pre-registration}: Full protocol registered on
    OSF or AsPredicted before data collection.
    \item \textbf{Time-locking}: Dream reports must be recorded,
    time-stamped, and locked before the future interval opens.
    \item \textbf{Polysomnographic staging}: Sleep-stage labels must be
    assigned from preregistered scoring rules by raters blind to dream
    content and outcome.
    \item \textbf{Double-blinding}: Experimenters interacting with
    subjects, dream judges, and data analysts must be blind to
    condition, stimulus schedule, and target identity until lock.
    \item \textbf{Matched decoys}: Each dream must be scored against
    one true future-event packet and decoy packets matched for
    category and base rate.
    \item \textbf{Judge reliability}: Inter-rater reliability on the
    primary dream-scoring rubric should satisfy
    Fleiss' $\kappa \geq 0.70$ before confirmatory unblinding.
    \item \textbf{Hardware RNG}: Future events determined by quantum
    or thermal random number generators with verified entropy.
    \item \textbf{Fixed primary endpoints}: Primary outcome measures
    (top-rank accuracy, mean reciprocal rank, regression slope,
    latency window) must be designated in advance.
    \item \textbf{Multiple-comparison correction}: Bonferroni or
    FDR (Benjamini-Hochberg) for multiple physiological channels.
    \item \textbf{Bayesian analysis}: Report both frequentist
    $p$-values and Bayes factors ($BF_{10}$) to distinguish
    ``evidence for null'' from ``inconclusive.''
    \item \textbf{Permutation controls}: Shuffle time labels and
    packet identities to generate null distributions.
\end{enumerate}

These standards are drawn from the methodological critiques of
Wagenmakers et al.\ \cite{Wagenmakers2012}, who argued that prior
presentiment studies suffered from optional stopping, flexible
analysis, and insufficient reporting of null results.


%=======================================================================
\section{Related Work}\label{sec:related}
%=======================================================================

\subsection{Two-Time Boundary Conditions in Physics}

The idea that future boundary conditions constrain present dynamics
has a long history. Hamilton's Principle of Least Action selects
trajectories by minimizing the action integral over a fixed time
interval with fixed endpoints \cite{Lanczos1970}. Feynman's path
integral sums amplitudes over all paths connecting initial and final
states. Most directly relevant is Aharonov's TSVF
\cite{Aharonov1964,Aharonov2008}, which describes quantum systems
using both a forward-evolving state (from preparation) and a
backward-evolving state (from post-selection). The present model is
a discrete, finite-dimensional analog: the window neutrality
constraint plays the role of post-selection, and the phantom
magnitude quantifies how strongly the post-selection constrains
the present.

The Wheeler-Feynman absorber theory \cite{WheelerFeynman1945}
similarly requires advanced (backward-in-time) solutions for
consistency, though these are usually interpreted as canceling in
the classical limit. Cramer's Transactional Interpretation
\cite{Cramer1986} extends this to a full interpretation of quantum
mechanics.

\subsection{Consciousness and Time}

The relationship between consciousness and temporal experience has
been explored from multiple directions. Tononi's Integrated
Information Theory (IIT) \cite{Tononi2004} proposes that
consciousness corresponds to integrated information ($\Phi$), which
depends on the causal structure of a system over time. The ``temporal
depth'' of $\Phi$ in IIT bears a structural resemblance to our
phantom magnitude: both quantify how much the past-future structure
of a system matters for its present state.

Global Workspace Theory (GWT) \cite{Baars1988,Dehaene2011} models
consciousness as a broadcast mechanism. In GWT terms, phantom light
would correspond to a constraint-derived signal that becomes
available for global broadcast during reduced-coherence states
(sleep), when the usual sensory competition for the workspace is
diminished.

\subsection{Dreams and Precognition}

The claim that dreams sometimes contain future information has a
long history in both anecdotal reports and experimental
investigation. Rhine \cite{Rhine1962} and Van de Castle
\cite{VanDeCastle1994} compiled systematic case collections.
Kramer \cite{Kramer2007} reviewed the dream-cognition literature
from a psychiatric perspective.

Hobson's AIM model \cite{Hobson2000,Hobson2005} explains dream
phenomenology through three dimensions: Activation, Input gating,
and chemical Modulation. In AIM terms, our coherence budget $C$
corresponds to the Activation dimension, and the prediction that
deep-sleep dreams carry stronger phantom content maps to the
low-A region of AIM state space.

\subsection{The Presentiment Literature}

Presentiment (physiological anticipation of random future stimuli)
has been reported by Radin \cite{Radin2004}, Bierman and Radin
\cite{Bierman1997}, McCraty et al.\ \cite{McCraty2004}, and
reviewed by Mossbridge et al.\ \cite{Mossbridge2012}, who conducted
a meta-analysis of 26 studies and reported a small but statistically
significant anticipatory effect.

However, this literature has been subject to serious methodological
critique. Wagenmakers et al.\ \cite{Wagenmakers2012} argued that
Bem's precognition results were artifacts of optional stopping and
flexible analysis. Alcock \cite{Alcock2003} reviewed the broader
history of parapsychological claims and argued that replication
failures are systematic. Rouder and Morey \cite{Rouder2011}
reanalyzed Bem's data using Bayesian methods and found the evidence
equivocal.

Our model does not assume the presentiment literature is correct.
It provides a framework within which presentiment, \emph{if it
exists}, would have specific properties (sleep-stage hierarchy,
debt scaling, response time law) that are not predicted by artifact
explanations. The falsification criteria in \S\ref{sec:predictions}
are designed to distinguish genuine phantom sensing from the
methodological artifacts identified by critics.


%=======================================================================
\section{Claim Hygiene}\label{sec:hygiene}
%=======================================================================

Every claim in this paper is classified below. ``PROVED'' means
machine-verified from the RS axioms together with any explicitly
declared modeling assumptions used in the statement. ``MODEL'' means
a definition, operationalization, or parameter choice that is
consistent with empirical data but not itself derived from axioms.
``HYPOTHESIS''
means a falsifiable empirical claim.

\begin{center}
\small
\begin{longtable}{p{4.2cm}cp{3.8cm}}
\toprule
\textbf{Claim} & \textbf{Status} & \textbf{Reference} \\
\midrule
\endhead
\multicolumn{3}{l}{\textit{RS Axioms (taken as given)}} \\
$J$ unique (T5) & AXIOM & \cite{RSCost2025} \\
$\varphi$ forced (T6) & AXIOM & \cite{RSAxioms2025} \\
$D=3$ (T8) & AXIOM & \cite{RSGeometry2025} \\
8-tick period (T7) & AXIOM & \cite{RSAxioms2025} \\
Window neutrality & AXIOM & \cite{RSAxioms2025} \\
Cost minimization & AXIOM & \cite{RSAxioms2025} \\
\midrule
\multicolumn{3}{l}{\textit{Phantom Light (proved from axioms)}} \\
Recognition forces balance & PROVED & Thm.~\ref{thm:lock} \\
$\Phi_{\mathrm{mag}} \geq 0$ & PROVED & Thm.~\ref{thm:phantom-props} \\
Monotone in $|\mathcal{D}|$ & PROVED & Thm.~\ref{thm:phantom-props} \\
Anti-monotone in $R$ & PROVED & Thm.~\ref{thm:phantom-props} \\
Zero debt $=$ zero phantom & PROVED & Thm.~\ref{thm:phantom-props} \\
Backward projection & PROVED & Thm.~\ref{thm:backward} \\
Cost inflation & PROVED & Thm.~\ref{thm:inflates} \\
Boundary reset & PROVED & Thm.~\ref{thm:reset} \\
\midrule
\multicolumn{3}{l}{\textit{Sleep model (proved from axioms + ordering)}} \\
Coherence ordering & MODEL & Def.~\ref{def:coherence} \\
Canonical integer budget & MODEL & Rem.~after Thm.~\ref{thm:stage3} \\
Stage window durations & MODEL & empirical \\
N3 $\geq$ N1 visibility & PROVED & Thm.~\ref{thm:stage3} \\
N3 $\geq$ REM visibility & PROVED & Thm.~\ref{thm:stage3} \\
Canonical budget: N3 $= 2\times$ N1 & PROVED & Thm.~\ref{thm:stage3} \\
N3 exposure dominates & PROVED & Thm.~\ref{thm:exposure} \\
Zero debt eliminates sleep phantom & PROVED & Thm.~\ref{thm:dreamless} \\
16 ultradians/day & PROVED & Thm.~\ref{thm:ultradians} \\
\midrule
\multicolumn{3}{l}{\textit{Consciousness interface}} \\
Sensitivity parameter & MODEL & \S\ref{sec:consciousness} \\
Readout positivity & PROVED & Thm.~\ref{thm:sensing} \\
Response time law & PROVED & Thm.~\ref{thm:response} \\
$\Theta$-coupling & PROVED & Thm.~\ref{thm:theta} \\
GCIC assumption & HYPOTHESIS & conditional \\
\midrule
\multicolumn{3}{l}{\textit{Empirical interface}} \\
Dream future-event proxy $\widehat{\mathcal{D}}_n(H)$ & MODEL & \S\ref{sec:operational} \\
Presentiment proxy $\widehat{\Phi}_t$ & MODEL & \S\ref{sec:operational} \\
Small-effect power envelope & MODEL & \S\ref{sec:effects} \\
\midrule
\multicolumn{3}{l}{\textit{Empirical predictions}} \\
Deep-sleep dominance & HYPOTHESIS & Pred.~\ref{pred:sleep} \\
Low debt, low match rate & HYPOTHESIS & Pred.~\ref{pred:debt} \\
Debt-content correlation & HYPOTHESIS & Pred.~\ref{pred:content} \\
Response time scaling & HYPOTHESIS & Pred.~\ref{pred:response} \\
Boundary reset periodicity & HYPOTHESIS & Pred.~\ref{pred:reset} \\
$\sigma > 0$ for humans & HYPOTHESIS & open \\
\bottomrule
\end{longtable}
\end{center}


%=======================================================================
\section{Discussion}\label{sec:discussion}
%=======================================================================

\subsection{What This Paper Does and Does Not Claim}

This paper presents a formal model, not a theory of precognition.
The model shows that within a specific mathematical framework (RS),
future conservation constraints create a computable cost inflation
in the present. The model then shows that this inflation has a
sleep-stage structure that matches the empirical coherence ordering.
Whether biological neural systems can detect this inflation is an
open empirical question.

The model does not require retrocausation, new particles, or
violation of any established physical law. It requires only the
RS axioms (which have independent motivation from information
theory and discrete dynamics \cite{RSAxioms2025,RSdalembert2026})
and the standard sleep-science finding that cortical coherence
decreases from waking through deep sleep.

It also does not claim that raw autobiographical narrative can be
read off directly from the latent RS variables. The operational
quantities in \S\ref{sec:operational} are empirical proxies intended
to preserve the model's monotonic structure under blinded,
preregistered scoring.

\subsection{Relationship to Existing Theories}

The phantom light mechanism is compatible with, but distinct from,
several existing theoretical frameworks:

\begin{itemize}
    \item \textbf{TSVF}: Both involve future boundary conditions
    constraining present states. Phantom Light provides a discrete,
    computable version with explicit sleep-stage predictions.
    \item \textbf{IIT}: Both define a quantity ($\Phi$ in IIT,
    $\Phi_{\mathrm{mag}}$ here) that measures the significance of
    temporal structure. IIT's $\Phi$ measures integrated information;
    Phantom Light's $\Phi_{\mathrm{mag}}$ measures outstanding
    balance debt.
    \item \textbf{GWT}: During sleep, reduced sensory competition
    for the global workspace could allow phantom signals to reach
    broadcast---consistent with our coherence-budget model.
\end{itemize}

\subsection{Limitations}

\begin{enumerate}
    \item The RS axioms are assumed, not derived in this paper. A
    reader who rejects the RS framework will reject the premises.
    The axioms are falsifiable independently
    \cite{RSAxioms2025,RSFermionMasses2025}.
    \item The coherence budget is a modeling choice. The qualitative
    predictions (ordering) are robust, but the exact ratios depend
    on the chosen integers.
    \item The empirical proxies
    $\widehat{\mathcal{D}}_n(H)$ and $\widehat{\Phi}_t$ are not the
    latent RS variables themselves. Their adequacy must be judged by
    prospective predictive performance under blinded protocols.
    \item The sensitivity parameter $\sigma$ is a free parameter at
    the modeling level. Its positivity and population scale are core
    empirical questions.
    \item The experimental predictions concern small effects in a
    domain with a troubled replication history. Stringent
    pre-registration and blinding are essential.
\end{enumerate}


%=======================================================================
\section{Conclusion}\label{sec:conclusion}
%=======================================================================

We have shown that within the RS framework, window neutrality creates
a two-time boundary condition in which recognition events generate
balance debts that the future must discharge. This balance debt
inflates the present cost function, creating a ``phantom
light''---a computable, bounded shadow of the future on the present
cost landscape. The mechanism requires no retrocausation.

A sleep-stage hierarchy emerges: lower cortical coherence increases
phantom visibility, with deep sleep (N3) dominating light sleep (N1)
and REM. Zero balance debt annihilates the effect in all modes.

All mathematical claims are machine-verified (52 theorems, zero
unproved placeholders). Every statement is classified as PROVED,
MODEL, or HYPOTHESIS. Falsification criteria with quantitative
thresholds are provided. Whether biological systems can read the
phantom field is an open question. We have therefore coupled the
formal structure to explicit operational proxies, a blinded
prospective dream-scoring design, and small-effect study-design
targets so that the empirical claims can be tested directly.


%=======================================================================
\begin{thebibliography}{99}
%=======================================================================

\bibitem{Aharonov1964}
Y.~Aharonov, P.~G.~Bergmann, and J.~L.~Lebowitz,
``Time symmetry in the quantum process of measurement,''
\textit{Phys.\ Rev.}\ \textbf{134}, B1410, 1964.

\bibitem{Aharonov2008}
Y.~Aharonov and L.~Vaidman,
``The Two-State Vector Formalism: An Updated Review,''
\textit{Lecture Notes in Physics} \textbf{734}, 399--447, 2008.

\bibitem{Alcock2003}
J.~E.~Alcock,
``Give the null hypothesis a chance: Reasons to remain doubtful
about the existence of psi,''
\textit{J.\ Consciousness Studies} \textbf{10}(6--7), 29--50, 2003.

\bibitem{Aserinsky1953}
E.~Aserinsky and N.~Kleitman,
``Regularly occurring periods of eye motility, and concomitant
phenomena, during sleep,''
\textit{Science} \textbf{118}, 273--274, 1953.

\bibitem{Baars1988}
B.~J.~Baars,
\textit{A Cognitive Theory of Consciousness},
Cambridge University Press, 1988.

\bibitem{Bierman1997}
D.~J.~Bierman and D.~I.~Radin,
``Anomalous anticipatory response on randomized future
conditions,''
\textit{Percept.\ Mot.\ Skills} \textbf{84}, 689--690, 1997.

\bibitem{Cramer1986}
J.~G.~Cramer,
``The transactional interpretation of quantum mechanics,''
\textit{Rev.\ Mod.\ Phys.}\ \textbf{58}, 647--687, 1986.

\bibitem{Dehaene2011}
S.~Dehaene and J.-P.~Changeux,
``Experimental and theoretical approaches to conscious processing,''
\textit{Neuron} \textbf{70}, 200--227, 2011.

\bibitem{Hobson2000}
J.~A.~Hobson, E.~F.~Pace-Schott, and R.~Stickgold,
``Dreaming and the brain: Toward a cognitive neuroscience of
conscious states,''
\textit{Behavioral and Brain Sciences} \textbf{23}(6), 793--842,
2000.

\bibitem{Hobson2005}
J.~A.~Hobson,
``Sleep is of the brain, by the brain and for the brain,''
\textit{Nature} \textbf{437}, 1254--1256, 2005.

\bibitem{Kleitman1963}
N.~Kleitman,
\textit{Sleep and Wakefulness},
University of Chicago Press, revised ed., 1963.

\bibitem{Kramer2007}
M.~Kramer,
\textit{The Dream Experience: A Systematic Exploration},
Routledge, 2007.

\bibitem{Lanczos1970}
C.~Lanczos,
\textit{The Variational Principles of Mechanics},
Dover Publications, 1970.

\bibitem{Massimini2005}
M.~Massimini et~al.,
``Breakdown of cortical effective connectivity during sleep,''
\textit{Science} \textbf{309}, 2228--2232, 2005.

\bibitem{McCraty2004}
R.~McCraty, M.~Atkinson, and R.~T.~Bradley,
``Electrophysiological evidence of intuition: Part~1. The
surprising role of the heart,''
\textit{J.\ Altern.\ Complement.\ Med.}\ \textbf{10}(1),
133--143, 2004.

\bibitem{Mossbridge2012}
J.~A.~Mossbridge, P.~Tressoldi, and J.~Utts,
``Predictive physiological anticipation preceding seemingly
unpredictable stimuli: A meta-analysis,''
\textit{Frontiers in Psychology} \textbf{3}, 390, 2012.

\bibitem{Radin2004}
D.~I.~Radin,
``Electrodermal presentiments of future emotions,''
\textit{J.\ Sci.\ Explor.}\ \textbf{18}, 253--274, 2004.

\bibitem{Rhine1962}
L.~E.~Rhine,
``Psychological processes in ESP experiences. Part I. Waking
experiences,''
\textit{J.\ Parapsychology} \textbf{26}, 88--111, 1962.

\bibitem{Rouder2011}
J.~N.~Rouder and R.~D.~Morey,
``A Bayes factor meta-analysis of Bem's ESP claim,''
\textit{Psychonomic Bulletin \& Review} \textbf{18}(4),
682--689, 2011.

\bibitem{RSAxioms2025}
J.~Washburn,
``The Algebra of Reality: A Recognition Science Derivation of
Physical Law,''
\textit{Axioms (MDPI)} \textbf{15}(2), 90, 2025.

\bibitem{RSCost2025}
J.~Washburn,
``Uniqueness of the Canonical Reciprocal Cost,''
\textit{Mathematics (MDPI)} \textbf{14}(6), 935, 2025.

\bibitem{RSdalembert2026}
J.~Washburn,
``The d'Alembert Inevitability Theorem,''
arXiv:2603.16237, 2026.

\bibitem{RSFermionMasses2025}
J.~Washburn and E.~Allahyarov,
``A Complete Derivation of the Fermion Spectrum from the
Recognition Composition Law,''
arXiv:2506.12859, 2025.

\bibitem{RSGeometry2025}
J.~Washburn,
``Recognition Geometry,''
\textit{Axioms (MDPI)} \textbf{15}(2), 2025.

\bibitem{RSLeanRepo}
J.~Washburn,
``Recognition Science Lean Repository,''
\url{https://github.com/jonwashburn/recognition-science}, 2026.

\bibitem{Tononi2004}
G.~Tononi,
``An information integration theory of consciousness,''
\textit{BMC Neuroscience} \textbf{5}, 42, 2004.

\bibitem{Tononi2006}
G.~Tononi,
``Consciousness, information integration, and the brain,''
\textit{Prog.\ Brain Res.}\ \textbf{150}, 109--126, 2006.

\bibitem{VanDeCastle1994}
R.~L.~Van de Castle,
\textit{Our Dreaming Mind},
Ballantine Books, 1994.

\bibitem{Wagenmakers2012}
E.-J.~Wagenmakers, R.~Wetzels, D.~Borsboom, and
H.~L.~J.~van der Maas,
``Why psychologists must change the way they analyze their data:
The case of psi,''
\textit{J.\ Personality and Social Psychology} \textbf{100}(3),
426--432, 2011.

\bibitem{Walker2017}
M.~P.~Walker,
\textit{Why We Sleep: Unlocking the Power of Sleep and Dreams},
Scribner, 2017.

\bibitem{WheelerFeynman1945}
J.~A.~Wheeler and R.~P.~Feynman,
``Interaction with the Absorber as the Mechanism of Radiation,''
\textit{Rev.\ Mod.\ Phys.}\ \textbf{17}(2--3), 157, 1945.

\end{thebibliography}


%=======================================================================
\appendix
%=======================================================================

\section{Lean Artifact Summary}\label{app:lean}

All theorems are verified in four Lean~4 modules:

\begin{center}
\small
\begin{tabular}{lrl}
\toprule
\textbf{Module} & \textbf{Theorems} & \textbf{Status} \\
\midrule
\texttt{PhantomLight.lean} & 17 & 0 sorry, 0 axiom \\
\texttt{PhantomSleep.lean} & 22 & 0 sorry, 0 axiom \\
\texttt{SleepCycles.lean} & 8 & 0 sorry, 0 axiom \\
\texttt{CayceBridge.lean} & 5 & 0 sorry, 0 axiom \\
\midrule
\textbf{Total} & \textbf{52} & \textbf{0 sorry, 0 axiom} \\
\bottomrule
\end{tabular}
\end{center}

To verify:
\begin{verbatim}
lake build IndisputableMonolith.Consciousness.PhantomSleep
\end{verbatim}
This transitively builds all dependencies. Source code is available
at \cite{RSLeanRepo}.


\end{document}
